\documentclass[CJK,13pt]{beamer}
\usepackage{CJKutf8}
\usepackage{beamerthemesplit}
\usetheme{Warsaw}
\useoutertheme[footline=authortitle]{miniframes}
\usepackage{amsmath}
\usepackage{amssymb}
\usepackage{graphicx}
\usepackage{eufrak}
\usepackage{color}
\usepackage{slashed}
\usepackage{simplewick}
\usepackage{tikz}
\usepackage{tcolorbox}
\usepackage{ulem}
%%figures
\def\lfig#1#2{\includegraphics[width=#1 in]{#2}}
\def\tfig#1#2{\includegraphics[height=#1 in]{#2}}
\def\addfig#1#2{\begin{center}\includegraphics[width=#1 in]{#2}\end{center}}
\def\question{\includegraphics[width=0.3in]{why.jpg}\,}
\def\answer{\includegraphics[width=0.3in]{baozou_haha.png}\,}
\def\wulian{\includegraphics[width=0.18in]{emoji_wulian.jpg}}
\def\bigwulian{\includegraphics[width=0.35in]{emoji_wulian.jpg}}
\def\bye{\includegraphics[width=0.18in]{emoji_bye.jpg}}
\def\bigbye{\includegraphics[width=0.35in]{emoji_bye.jpg}}
\def\huaixiao{\includegraphics[width=0.18in]{emoji_huaixiao.jpg}}
\def\bighuaixiao{\includegraphics[width=0.35in]{emoji_huaixiao.jpg}}
\def\jianxiao{\includegraphics[width=0.18in]{emoji_jianxiao.jpg}}
\def\bigjianxiao{\includegraphics[width=0.35in]{emoji_jianxiao.jpg}}
\def\haoqi{\includegraphics[width=0.18in]{emoji_haoqi.jpg}}
%% colors
\def\blacktext#1{{\color{black}#1}}
\def\bluetext#1{{\color{blue}#1}}
\def\redtext#1{{\color{red}#1}}
\def\darkbluetext#1{{\color[rgb]{0,0.2,0.6}#1}}
\def\skybluetext#1{{\color[rgb]{0.2,0.7,1.}#1}}
\def\cyantext#1{{\color[rgb]{0.,0.5,0.5}#1}}
\def\greentext#1{{\color[rgb]{0,0.7,0.1}#1}}
\def\darkgray{\color[rgb]{0.2,0.2,0.2}}
\def\lightgray{\color[rgb]{0.6,0.6,0.6}}
\def\gray{\color[rgb]{0.4,0.4,0.4}}
\def\blue{\color{blue}}
\def\red{\color{red}}
\def\orange{\color[rgb]{1.,0.8,0.}}
\def\green{\color{green}}
\def\darkgreen{\color[rgb]{0,0.4,0.1}}
\def\darkblue{\color[rgb]{0,0.2,0.6}}
\def\skyblue{\color[rgb]{0.2,0.7,1.}}
%%control
\def\ulim#1{\underset{#1}{\lim}}
\def\heaviside{\,\mathrm{h}\,}
\def\bral{\left(\begin{array}{l}}
\def\brar{\end{array}\right)}
\def\brall{\left(\begin{array}{ll}}
\def\brarr{\end{array}\right)}
\def\bralll{\left(\begin{array}{lll}}
\def\brarrr{\end{array}\right)}
\def\branchl{\left\{\begin{array}{l}}
\def\branchr{\end{array}\right.}
\def\branchll{\left\{\begin{array}{ll}}
\def\branchrr{\end{array}\right.}
\def\branchlll{\left\{\begin{array}{lll}}
\def\branchrrr{\end{array}\right.}
\def\sfgamma#1{\,\Gamma\left( #1 \right)\,}
\def\be{\begin{equation}}
\def\ee{\nonumber\end{equation}}
\def\bea{\begin{eqnarray}}
\def\eea{\nonumber\end{eqnarray}}
\def\bch{\begin{CJK}{UTF8}{gbsn}}
\def\ech{\end{CJK}}
\def\bitem{\begin{itemize}}
\def\eitem{\end{itemize}}
\def\bcenter{\begin{center}}
\def\ecenter{\end{center}}
\def\bex{\begin{minipage}{0.2\textwidth}\includegraphics[width=0.6in]{jugelizi.png}\end{minipage}\begin{minipage}{0.76\textwidth}}
\def\eex{\end{minipage}}
\def\chtitle#1{\frametitle{\bch#1\ech}}
\def\bmat#1{\left(\begin{array}{#1}}
\def\emat{\end{array}\right)}
\def\bcase#1{\left\{\begin{array}{#1}}
\def\ecase{\end{array}\right.}
\def\bmini#1{\begin{minipage}{#1\textwidth}}
\def\emini{\end{minipage}}
\def\tbox#1{\begin{tcolorbox}#1\end{tcolorbox}}
\def\pfrac#1#2#3{\left(\frac{\partial #1}{\partial #2}\right)_{#3}}
\def\res#1#2{\,\mathrm{res}\,#1\left(#2\right)\,}
\def\newt#1#2{\left(\begin{array}{c}#1\\ #2\end{array}\right)}
\def\reof#1{\,\mathrm{Re}{\left(#1\right)\,}}
\def\imof#1{\,\mathrm{Im}{\left(#1\right)\,}}
\def\Arg#1{\,\mathrm{Arg}\,#1\,}
%%symbols
\def\intfull{\int_{-\infty}^\infty}
\def\inthalf{\int_0^\infty}
\def\bropt{\,(\ \ \ )}
\def\sech{\mathrm{sech}\,}
\def\csch{\mathrm{csch}\,}
\def\sone{$\star$}
\def\stwo{$\star\star$}
\def\sthree{$\star\star\star$}
\def\sfour{$\star\star\star\star$}
\def\sfive{$\star\star\star\star\star$}
\def\rint{{\int_\leftrightarrow}}
\def\roint{{\oint_\leftrightarrow}}
\def\stdHf{{\textit{\r H}_f}}
\def\deltaH{{\Delta \textit{\r H}}}
\def\ii{{\dot{\imath}}}
\def\skipline{{\vskip0.1in}}
\def\skiplines{{\vskip0.2in}}
\def\lagr{{\mathcal{L}}}
\def\hamil{{\mathcal{H}}}
\def\vecv{{\mathbf{v}}}
\def\vecf{{\mathbf{f}}}
\def\vecg{{\mathbf{g}}}
\def\vecupsilon{{\mathbf{\upsilon}}}
\def\vecu{{\mathbf{u}}}
\def\vecj{{\mathbf{j}}}
\def\vecx{{\mathbf{x}}}
\def\vecy{{\mathbf{y}}}
\def\vecz{{\mathbf{z}}}
\def\veck{{\mathbf{k}}}
\def\vecp{{\mathbf{p}}}
\def\vecn{{\mathbf{n}}}
\def\vecA{{\mathbf{A}}}
\def\vecP{{\mathbf{P}}}
\def\vecsigma{{\mathbf{\sigma}}}
\def\hatJn{{\hat{J_\vecn}}}
\def\hatJx{{\hat{J_x}}}
\def\hatJy{{\hat{J_y}}}
\def\hatJz{{\hat{J_z}}}
\def\hatj#1{\hat{J_{#1}}}
\def\hatphi{{\hat{\phi}}}
\def\hatq{{\hat{q}}}
\def\hatpi{{\hat{\pi}}}
\def\vel{\upsilon}
\def\Dint{{\mathcal{D}}}
\def\adag{{\hat{a}^\dagger}}
\def\bdag{{\hat{b}^\dagger}}
\def\cdag{{\hat{c}^\dagger}}
\def\ddag{{\hat{d}^\dagger}}
\def\hata{{\hat{a}}}
\def\hatb{{\hat{b}}}
\def\hatc{{\hat{c}}}
\def\hatd{{\hat{d}}}
\def\hatD{{\,\hat{D}}}
\def\hatN{{\hat{N}}}
\def\hatH{{\hat{H}}}
\def\hatp{{\hat{p}}}
\def\Fup{{F^{\mu\nu}}}
\def\Fdown{{F_{\mu\nu}}}
\def\newl{\nonumber \\}
\def\vece{\mathrm{e}}
\def\calM{{\mathcal{M}}}
\def\calT{{\mathcal{T}}}
\def\calR{{\mathcal{R}}}
\def\barpsi{\bar{\psi}}
\def\baru{\bar{u}}
\def\barv{\bar{\upsilon}}
\def\qeq{\stackrel{?}{=}}
\def\ftf{\stackrel{\mathcal{FT}}{\Longrightarrow}}
\def\ftb{\stackrel{\mathcal{FT}}{\Longleftarrow}}
\def\ftfb{\stackrel{\mathcal{FT}}{\Longleftrightarrow}}
\def\ltf{\stackrel{\mathcal{LT}}{\Longrightarrow}}
\def\ltb{\stackrel{\mathcal{LT}}{\Longleftarrow}}
\def\ltfb{\stackrel{\mathcal{LT}}{\Longleftrightarrow}}
\def\torder#1{\mathcal{T}\left(#1\right)}
\def\rorder#1{\mathcal{R}\left(#1\right)}
\def\contr#1#2{\contraction{}{#1}{}{#2}#1#2}
\def\trof#1{\mathrm{Tr}\left(#1\right)}
\def\trace{\mathrm{Tr}}
\def\comm#1{\ \ \ \left(\mathrm{used}\ #1\right)}
\def\tcomm#1{\ \ \ (\text{#1})}
\def\slp{\slashed{p}}
\def\slk{\slashed{k}}
\def\calp{{\mathfrak{p}}}
\def\veccalp{\mathbf{\mathfrak{p}}}
\def\Tthree{T_{\tiny \textcircled{3}}}
\def\pthree{p_{\tiny \textcircled{3}}}
\def\dbar{{\,\mathchar'26\mkern-12mu d}}
\def\erf{\mathrm{erf}}
\def\const{\mathrm{const.}}
\def\pheat{\pfrac p{\ln T}V}
\def\vheat{\pfrac V{\ln T}p}
%%units
\def\fdeg{{^\circ \mathrm{F}}}
\def\cdeg{^\circ \mathrm{C}}
\def\atm{\,\mathrm{atm}}
\def\angstrom{\,\text{\AA}}
\def\SIL{\,\mathrm{L}}
\def\SIkm{\,\mathrm{km}}
\def\SIyr{\,\mathrm{yr}}
\def\SIGyr{\,\mathrm{Gyr}}
\def\SIV{\,\mathrm{V}}
\def\SImV{\,\mathrm{mV}}
\def\SIeV{\,\mathrm{eV}}
\def\SIkeV{\,\mathrm{keV}}
\def\SIMeV{\,\mathrm{MeV}}
\def\SIGeV{\,\mathrm{GeV}}
\def\SIcal{\,\mathrm{cal}}
\def\SIkcal{\,\mathrm{kcal}}
\def\SImol{\,\mathrm{mol}}
\def\SIN{\,\mathrm{N}}
\def\SIHz{\,\mathrm{Hz}}
\def\SIm{\,\mathrm{m}}
\def\SIcm{\,\mathrm{cm}}
\def\SIfm{\,\mathrm{fm}}
\def\SImm{\,\mathrm{mm}}
\def\SInm{\,\mathrm{nm}}
\def\SImum{\,\mathrm{\mu m}}
\def\SIJ{\,\mathrm{J}}
\def\SIW{\,\mathrm{W}}
\def\SIkJ{\,\mathrm{kJ}}
\def\SIs{\,\mathrm{s}}
\def\SIkg{\,\mathrm{kg}}
\def\SIg{\,\mathrm{g}}
\def\SIK{\,\mathrm{K}}
\def\SImmHg{\,\mathrm{mmHg}}
\def\SIPa{\,\mathrm{Pa}}
%page
\def\secpage#1#2{\begin{frame}\bcenter{\bf \Huge #1} \ecenter \skipline \tbox{\bcenter #2 \ecenter}\end{frame}}
\def\append#1#2{\secpage{附录#1}{#2}}
\def\thinka#1{\begin{frame}\frametitle{思考题}\bcenter\lfig{0.4}{think.jpg}  \skipline #1 \ecenter\end{frame}}
\def\thinkb#1{\begin{frame}\frametitle{思考题}\bcenter\lfig{0.5}{think1.jpg}  \skipline #1 \ecenter\end{frame}}
\def\thinkc#1{\begin{frame}\frametitle{思考题}\bcenter\lfig{0.5}{think2.jpg}  \skipline #1 \ecenter\end{frame}}
\def\thinkd#1{\begin{frame}\frametitle{思考题}\bcenter\lfig{0.5}{think3.jpg}  \skipline #1 \ecenter\end{frame}}
\def\thinke#1{\begin{frame}\frametitle{思考题}\bcenter\lfig{0.5}{think4.jpg}  \skipline #1 \ecenter\end{frame}}

\def\courseurl{http://zhiqihuang.top/cosm}

\def\tpage#1#2{
\title{COSM \S{#1}  #2}
  \author{Zhiqi Huang}

\begin{frame}
\begin{center}
{\bf \Large Cosmology}

{\vskip 0.1in}

{\Large \S #1 #2}

{\vskip 0.2in}

{Lecturer: 黄志琦}

\vskip 0.2in

\courseurl

\end{center}
\end{frame}
}


  \date{}
  \begin{document}
  \bch
  \tpage{2}{FRW metric}

  \begin{frame}
    \frametitle{Mach的思想实验}

    \addfig{1.5}{mach_experiment.jpg}

    当水桶相对宇宙的恒星/星系背景没有旋转时，离心力就消失了，这是惊人的巧合吗？
  \end{frame}


  \begin{frame}
    \frametitle{标准宇宙学的基本图像}
    \bmini{0.42}
    宇宙由一个非常均匀且各向同性的初始状态开始，相距很远的“测地线粒子”间的距离用标准尺子——普朗克长度来度量，不断地增大（即所谓的宇宙大爆炸）。不过，邻近的粒子间由于引力的不稳定性，逐渐结团。
    \emini
    \bmini{0.53}
    \addfig{2}{large_scale_structure.jpg}
    \emini

    \skiplines

    虽然引力导致物质在小尺度上结团，在亿光年为单位的大尺度上看，宇宙仍然均匀和各向同性。    

  \end{frame}

\begin{frame}
  \frametitle{宇宙学基本原理(cosmological principle)}
  从天圆地方，到地心说，再到日心说，再到发现河外星系……一次又一次的沉重打击之后，人类基本上已经放弃了把自己置于宇宙中心的想法。

  \addfig{2}{universe.jpg}
  

  
  \skipline
  
  宇宙学基本原理：\sout{既然我们做不了宇宙中心，那谁也别想做。}{\blue 宇宙中没有任何特殊观测者，大尺度上看，宇宙均匀且各向同性。}
  {\scriptsize (其实叫基本假设更合理)}
\end{frame}
  

\begin{frame}
  \frametitle{FRW 度规}
  满足均匀且各向同性的只有 Friedmann–Lema\^itre–Robertson–Walker 度规，有时候也简称 Friedmann–Robertson–Walker 或 FLRW 或 FRW 度规：

  {\scriptsize Lema\^itre 的失败可能在于英文键盘打不出他的名字……}
  
  $$ ds^2 = dt^2 - a^2(t)\left(\frac{dr^2}{1-kr^2}+r^2d\theta^2+r^2\sin^2\theta d\phi^2\right)$$
  这里的 $a(t)$ 是任意一个 $t$ 的光滑函数。$k$ 是代表着三维空间的曲率的常数。

  \skipline
  
  为了书写方便，在这一讲里我们作一些约定：
  \bitem
\item{用符号上一点表示对宇宙学时间$t$的导数：如 $\dot{a}$ 表示 $\frac{da}{dt}$。}
\item{用下标``0''表示一个参量的“目前”值。例如$t_0$ 是“目前”的宇宙学时间。}
  \eitem
\end{frame}

\begin{frame}
  \frametitle{共动坐标}
  由于 $\Gamma^r_{\ tt}=\Gamma^\theta_{\ tt} =\Gamma^{\phi}_{\ tt}=0$，如果一开始测试粒子的 $p^r=p^\theta=p^\phi = 0$，则根据测地线方程
  $$\frac{d p^\alpha}{dt} + \Gamma^{\alpha}_{\mu\nu}\frac{p^\mu p^\nu}{p^0} = 0$$
  这里的 $\alpha=r, \theta,\phi$，无论 $\mu,\nu$ 同时为 $t$ (这样 $\Gamma^{\alpha}_{\mu\nu}$ 为零），还是 $\mu,\nu$ 至少有一个不是 $t$（这样 $p^\mu p^\nu$ 为零)，都会使 $\frac{d p^\alpha}{dt}$ 为零。

      \addfig{2}{chuangmoli.jpg}
      
  也就是说 ，测地线粒子在坐标 $(r,\theta,\phi)$ 可以一直保持“静止不动”，我们把 $r,\theta,\phi$ 叫做“共动坐标”。  
\end{frame}


\begin{frame}
  \frametitle{尺度因子 $a$ 的归一化是任意的} 
  FRW度规里的 $a(t)$ 称为{\blue 宇宙学尺度因子(scale factor)}。

  \skipline
  
  $a(t)r$ 具有物理长度的意义。如果把 $a(t)$ 多乘一个常数，同时把 $r$ 除以这个常数，并把 $k$ 乘以这个常数的平方，则不改变任何物理内容。我们经常把 $a_0$ (今天的 $a$)取成 $1$。
  
\end{frame}


\begin{frame}
  \frametitle{空间曲率 $k$}
  用固定的宇宙学时间 $t$ 在FRW时空切出一个三维空间，并把该时刻的 $a$ 归一化为 $1$，则三维空间的度规为：
  $$\left.ds^2\right\vert_{3D} = \frac{dr^2}{1-kr^2}+r^2d\theta^2+r^2\sin^2\theta d\phi^2.$$
  这时可以计算出 Ricci标量为 $6k$。

  \skiplines
  
  我们把 $k>0$ （正曲率空间）的宇宙称为闭合宇宙(closed universe)，$k<0$ 的宇宙叫做开放宇宙(open universe)，把 $k=0$ 的宇宙叫做平坦宇宙(flat universe)。

  {\scriptsize 注意“平坦”只是指空间平坦，整个四维时空还是弯曲的。}
  
\end{frame}

\begin{frame}
\frametitle{平坦宇宙和观测数据（特别是宇宙微波背景辐射数据）符合得还不错}
  \addfig{3}{spatial_curvature.jpg}  
\end{frame}

\begin{frame}
  \frametitle{哈勃参数}
  哈勃参数的定义为：
  \be
  H = \frac{\dot a}{a}
  \ee
  在标准宇宙学里，它恒大于零（即宇宙一直在膨胀）。

  \skipline

  $H_0$，也就是在 $t=t_0$ 时刻的 $H$，叫做{\blue 哈勃常数}。它大致在 $70\mathrm{km/s/Mpc}$ 左右。

  {\scriptsize 这里的长度单位 $\mathrm{Mpc}$ 是百万秒差距的意思，一个秒差距大约是 $3.26$ 光年。}

\end{frame}


\begin{frame}
  \frametitle{Hubble tension}
  \addfig{3.5}{H0tension.png}
\end{frame}

\begin{frame}
  \frametitle{思考题}
  \addfig{0.5}{think5.jpg}
  
  假设一个共动观测者在 $(t_1, r_1, \theta_1,\phi_1)$ 发射了一个TA看来频率为 $\nu$ 的光子。在 $(t_2, r_2,\theta_2, \phi_2)$ 处一个共动观测者接收到光子，接收者看到光子的频率是多少?
\end{frame}


\begin{frame}
  \frametitle{宇宙学红移}
  共动观测者的四维速度为 $u_\mu = (1, 0, 0, 0)$，因此看到的光子频率直接和光子在FRW坐标系的 $p^t$ 成正比。光子的测地线方程（见附录的联络表达式)为:
  $$ \frac{dp^t}{dt} + \frac{a\dot a}{p^t} \left[\frac{(p^r)^2}{1-kr^2} + r^2\left(p^\theta\right)^2 + r^2\sin^2\theta \left(p^\phi\right)^2\right] = 0.$$
    又光子走的是零测地线，所以：
    $$\left(p^t\right)^2 - a^2\left[\frac{(p^r)^2}{1-kr^2} + r^2\left(p^\theta\right)^2 + r^2\sin^2\theta \left(p^\phi\right)^2 \right] = 0.$$
    于是就有
    $$\frac{d p^t}{dt}  + \frac{\dot a}{a } p^t = 0$$
    它的解显然是 $p^t \propto \frac{1}{a} $。也就是说，光子的波长和 $a$ 成正比。
\end{frame}


\begin{frame}
  \frametitle{宇宙学红移}
  宇宙学红移就定义为波长的相对变化量
  $$z \equiv \frac{\lambda_{\rm obs}-\lambda_{\rm emit}}{\lambda_{\rm emit}}$$
  通常观测者指的都是目前地球上的观测者。所以在宇宙学时间 $t$ 发射出的光子的宇宙学红移为：
  $$ z = \frac{a_0-a}{a}, $$
  这里的 $a=a(t)$ 是发射点的尺度因子。 容易看出
  $$  \frac{a}{a_0} = \frac{1}{1+z}.$$
\end{frame}


\begin{frame}
  \frametitle{观测量的红移}
  在天文上，经常看到“这个星系的红移是$z$”这种说法。它的意思是：我们收集到的来自星系的光子，是星系在
  $$ \frac{a(t)}{a_0} = \frac{1}{1+z}$$
  对应的宇宙学时间 $t$ 时刻发出的(假设星系足够远，可以近似看成一个共动测试粒子)。

  \skipline

  显然，$z$ 越大，代表看到的星系越古老，离我们的距离也越远。
\end{frame}


\begin{frame}
  \frametitle{Friedmann Equations}
  如果宇宙的物质可以看成均匀的理想流体，密度为 $\rho(t)$, 压强为 $p(t)$，则爱因斯坦方程可以简化为：

  \tbox{
  \bea
  \frac{k + \dot a ^2}{a^2} &=& \frac{8\pi G}{3}\rho \newl
  \frac{\ddot a}{a} &=& -\frac{4\pi G}{3}\left(\rho + 3p\right)  
  \eea
  }
  这两个方程分别称为第一个和第二个 {\blue Friedmann 方程}。
  
\end{frame}


\begin{frame}
  \frametitle{附录1：FRW度规的联络}
\bmini{0.48}
\bea
\Gamma^t_{\ rr} &=& \frac{a\dot a}{1-kr^2} \newl
\Gamma^t_{\ \theta \theta } &=& r^2a\dot a \newl
\Gamma^t_{\ \phi \phi } &=& r^2a\dot a \sin^2\theta \newl
\Gamma^r_{\ r t} &=& \frac{\dot a}{a} \newl 
\Gamma^r_{\ rr} &=&  \frac{k r}{1-k r^2}  \newl
\Gamma^r_{\ \theta \theta } &=& -r(1-kr^2) \newl
\Gamma^r_{\ \phi \phi } &=& -r \left(1-k r^2 \right) \sin^2\theta
\eea
\emini
\bmini{0.48}
\bea
\Gamma^\theta _{\ \theta  t} &=& \frac{\dot a}{a} \newl
\Gamma^\theta _{\ \theta r} &=& \frac{1}{r} \newl
\Gamma^\theta _{\ \phi \phi } &=& - \sin\theta \cos\theta \newl
\Gamma^\phi _{\ \phi  t} &=& \frac{\dot a}{a} \newl
\Gamma^\phi _{\ \phi r} &=& \frac{1}{r} \newl
\Gamma^\phi _{\ \phi \theta } &=& \cot\theta
\eea
\emini
\end{frame}


\begin{frame}
  \frametitle{附录2：FRW度规的Ricci张量和Einstein 张量}
  Ricci tensor:
  \bea
R^{ t}_{\ t}&=& - \frac{3 \ddot a}{a} \newl
R^{r}_{\ r}&=&R^{\theta }_{\ \theta } =R^{\phi }_{\ \phi }= - \frac{2 k + a \ddot a + 2 \dot a^2}{a^2} 
\eea
Ricci scalar:
\be
R= - \frac{6\left(k + a \ddot a + \dot a^2\right)}{a^2} 
\ee
Einstein tensor:
\bea
G^{ t}_{\ t}&=& \frac{3\left(k + \dot a ^2\right)}{a^2}  \nonumber \\
G^{r}_{\ r}&=& G^{\theta }_{\ \theta }= G^{\phi }_{\ \phi } =\frac{k + 2 a \ddot a + \dot a ^2}{a^2}  
\eea

\end{frame}


    \ech
\end{document}




  
